% Generated draft: requires mathematical and source review.
\documentclass[11pt]{article}
\usepackage[margin=25mm]{geometry}
\usepackage{fontspec}
\setmainfont{Latin Modern Roman}
\usepackage{amsmath,amssymb,mathtools,unicode-math}
\setmathfont{Latin Modern Math}
\usepackage{enumitem,needspace,xurl,hyperref}
\setlength{\parindent}{0pt}
\setlength{\parskip}{5pt}
\setlength{\emergencystretch}{3em}
\newcommand{\N}{\mathbb{N}}
\newcommand{\Z}{\mathbb{Z}}
\newcommand{\Q}{\mathbb{Q}}
\newcommand{\R}{\mathbb{R}}
\newcommand{\C}{\mathbb{C}}
\newcommand{\F}{\mathbb{F}}
\newcommand{\A}{\mathbb{A}}
\newcommand{\PP}{\mathbb{P}}
\newcommand{\E}{\mathbb{E}}
\newcommand{\eps}{\varepsilon}
\newcommand{\dd}{\,\mathrm{d}}
\begin{document}
\section*{UnsolvedMath ID: 30001637}
\textbf{Source problem code:} catalogue entry
\subsection{Definitions and mathematical statement}
The supplied record leaves the surface conventions and the range and quantification of the genus unspecified. Under the conventional interpretation, let \(S_g\) be a closed, connected, orientable surface of integer genus \(g\), with no boundary or marked points, and let
\[
G_g=\operatorname{MCG}(S_g)
=\pi_0\bigl(\operatorname{Homeo}^{+}(S_g)\bigr)
\]
be its group of orientation-preserving homeomorphisms modulo isotopy. These surface and group conventions are conditional interpretations requiring confirmation from the source.

Interpret a cell complex as a CW complex. A space \(X\) is contractible if its identity map is homotopic to a constant map. For the discrete group \(G_g\), interpret a continuous action on \(X\) as proper if, for every compact subset \(K\subseteq X\), the set
\[
\{h\in G_g:hK\cap K\ne\varnothing\}
\]
is finite. The action is cocompact if \(G_g\backslash X\) is compact. The dimension of \(X\) is the supremum of its cell dimensions.

With these conventions, the existence question supplied by the record is: for an intended genus \(g\), does there exist a contractible cell complex \(X_g\), together with a proper cocompact action of \(G_g\), such that
\[
\dim X_g=4g-5?
\]
The dimension is stated as an equality. A nonempty CW complex has nonnegative dimension, so the conventional integer-genus formulation can only have an affirmative answer when \(g\ge 2\). The record does not specify whether the question concerns every such genus, a fixed genus, or a narrower range. It also does not specify whether the action must preserve the cell structure.

The additional question, whether there are obstructions to the existence of such a complex and action, is retained as an open-ended request: no particular obstruction, invariant, or necessary condition is specified. The supplied dated assessment describes the extracted statement as only the fragment ``Are there obstructions'', whereas the supplied statement fields contain the fuller existence question above. This discrepancy prevents treating the fuller formulation as independently established source wording.
\subsection{Short English statement}
The supplied question asks whether the mapping class group of a genus-\(g\) surface admits a proper cocompact action on a contractible cell complex of dimension exactly \(4g-5\), and what might obstruct such an action. The genus quantifier, surface conventions, action requirements, and intended meaning of obstruction require review; the supplied open status is dated curation.
\subsection{Sources}
\begin{itemize}
\item Supplied local UnsolvedMath registry record, permanent ID \(30001637\), source problem code \texttt{catalogue entry}, titled ``Minimal-Dimensional Cocompact Models for Mapping Class Groups'', in the Oberwolfach Reports Open Problems collection. The record supplies proposed year \(2010\) and names no proposer.
\item Supplied source citation: \emph{Teichmüller Theory} (2011), \emph{Oberwolfach Reports}, pp.~3085--3157, DOI \texttt{10.4171/owr/2010/53}. Source URL and the supplied literature-source link labelled ``original Oberwolfach problem report'': \url{https://doi.org/10.4171/owr/2010/53}. The record describes Oberwolfach Reports as documenting mathematical workshops and their open-problem discussions. No problem-specific page is supplied.
\item Supplied literature assessment, checked on \texttt{2026-08-21}: status ``open''. It says that the extracted statement is only the fragment ``Are there obstructions'', does not specify a theorem, and that the underlying cell-complex existence question remains unresolved in the cited report. This is supplied dated curation derived from the supplied OWR material, not independently verified here; time-sensitive status and cited claims require independent verification.
\item The registry labels the statement \texttt{corrected\_verified} and reports that its mathematical content was checked against the official Oberwolfach report, with notation, context, and question form normalized without changing the source-level claim. This verification assertion is supplied metadata, not a verification performed for this draft.
\end{itemize}
\end{document}
